%!TEX root = ../draft_main.tex
\section{Complementarity relation between chaometer's and CJ's purity}
We want the Haar average (over pure input states $\ket{\psi}$) of the output
purity for a single-qubit channel $\mathcal{E}$:
\begin{equation}
\mathbb{E}_{\psi\sim \mathrm{Haar}}\Big[
\Tr\big(\mathcal{E}(\ket{\psi}\bra{\psi})^2\big)
\Big].
\end{equation}

For any two operators $A,B$ on the same Hilbert space,
\begin{equation}
\Tr(AB) = \Tr\big((A \otimes B) S \big),
\end{equation}
where $S$ is the swap operator on the tensor product of two copies of the
output space. Therefore, for a pure input $\ket{\psi}$,
\begin{equation}
\begin{aligned}
\Tr\big(\mathcal{E}(\ket{\psi}\bra{\psi})^2\big) 
&= \Tr\Big( \big( \mathcal{E}(\ket{\psi}\bra{\psi}) \otimes 
\mathcal{E}(\ket{\psi}\bra{\psi}) \big) S \Big) \\
%
&= \Tr\Big( (\mathcal{E} \otimes \mathcal{E})
(\ket{\psi}\bra{\psi}^{\otimes 2}) S \Big).
\end{aligned}
\end{equation}

The Haar average of the two-copy projector is the projector onto the
symmetric subspace, normalized:
\begin{equation}
\int \ket{\psi}\bra{\psi}^{\otimes 2} \, d\psi
= \frac{P_{\mathrm{sym}}}{\dim P_{\mathrm{sym}}}
= \frac{I + F}{d(d+1)},
\end{equation}
where $F$ is the input-space swap operator and $d$ is the input dimension
(for a qubit, $d=2$). Thus,
\begin{equation}
\begin{aligned}
\mathbb{E}_\psi \big[ \Tr(\mathcal{E}(\ket{\psi}\bra{\psi})^2) \big]
&= \Tr \Big( (\mathcal{E} \otimes \mathcal{E}) 
\big( \frac{I+F}{d(d+1)} \big) S \Big) \\
&= \frac{1}{d(d+1)} \Big(
\Tr\big( (\mathcal{E} \otimes \mathcal{E})(I) \, S \big) \\
&\quad + \Tr\big( (\mathcal{E} \otimes \mathcal{E})(F) \, S \big)
\Big).
\end{aligned}
\end{equation}

For the first term,
\begin{equation}
(\mathcal{E} \otimes \mathcal{E})(I) = \mathcal{E}(I) \otimes \mathcal{E}(I),
\end{equation}
and using
\begin{equation}
\Tr(A \otimes B \, S) = \Tr(AB)
\end{equation}
gives
\begin{equation}
\Tr\big( (\mathcal{E} \otimes \mathcal{E})(I) \, S \big) = 
\Tr\big(\mathcal{E}(I)^2 \big).
\end{equation}

For the second term, expand with an input orthonormal basis $\{\ket{i}\}$:
\begin{equation}
F = \sum_{i,j} \ket{i}\bra{j} \otimes \ket{j}\bra{i},
\end{equation}
so
\begin{equation}
(\mathcal{E} \otimes \mathcal{E})(F) = 
\sum_{i,j} \mathcal{E}(\ket{i}\bra{j}) \otimes \mathcal{E}(\ket{j}\bra{i}),
\end{equation}
and again $\Tr(A \otimes B \, S) = \Tr(AB)$ gives
\begin{equation}
\Tr\big( (\mathcal{E} \otimes \mathcal{E})(F) \, S \big)
= \sum_{i,j} \Tr \big( \mathcal{E}(\ket{i}\bra{j})
\mathcal{E}(\ket{j}\bra{i}) \big).
\end{equation}

The right-hand side is exactly $\Tr(\mcD^2)$ if one uses the
(un-normalized) \jami{} matrix
\begin{equation}
\mcD = \sum_{i,j} \mathcal{E}(\ket{i}\bra{j}) \otimes
\ket{i}\bra{j},
\end{equation}
because a short calculation shows
\begin{equation}
\Tr(\mcD^2) = 
\sum_{i,j} \Tr(\mathcal{E}(\ket{i}\bra{j}) \mathcal{E}(\ket{j}\bra{i})).
\end{equation}

Hence, for general input dimension $d$,
\begin{equation}
\mathbb{E}_\psi \big[ \Tr(\mathcal{E}(\ket{\psi}\bra{\psi})^2) \big]
= \frac{1}{d(d+1)} \Big( \Tr[\mathcal{E}(I)^2] + \Tr[\mcD^2] \Big).
\end{equation}

For a qubit,
\begin{equation}
\mathbb{E}_\psi \big[ \Tr(\mathcal{E}(\ket{\psi}\bra{\psi})^2) \big]
= \frac{1}{6} \Big( \Tr[\mathcal{E}(I)^2] + \Tr[\mcD^2] \Big).
\end{equation}
Using a state 2-desing we can rewrite the above expression as
\begin{equation}
\sum_{k=1}^{6} \Tr\Big[ \mathcal{E}(\dyad{\phi_k})^2 \Big]
= 4\Tr[\mathcal{E}(I/2)^2] + \Tr[\mcD^2],
\end{equation}
where $\ket{\phi_k}$ are the six eigenstates of the three Pauli matrices.

\subsubsection*{Remarks}
\begin{itemize}
\item If $\mathcal{E}$ is trace-preserving, $\Tr[\mathcal{E}(I)] = \Tr[I] = d$,
  but $\Tr[\mathcal{E}(I)^2]$ depends on whether $\mathcal{E}(I)$ is the
  identity (unital) or not.
\item For a \textbf{unitary channel} $\mathcal{E}(\rho) = U \rho U^\dagger$,
  one finds $\mathcal{E}(I) = I$ and $\Tr[\mcD^2] = d^2$, so the
  formula returns $1$ as expected.
\item In Kraus form $\mathcal{E}(\rho) = \sum_\alpha K_\alpha \rho K_\alpha^\dagger$,
  one can write
\begin{equation}
\Tr[\mcD^2] = \sum_{i,j} \Tr \Big(
\sum_\alpha K_\alpha \ket{i}\bra{j} K_\alpha^\dagger
\sum_\beta K_\beta \ket{j}\bra{i} K_\beta^\dagger \Big),
\end{equation}
which is the explicit (but sometimes less compact) Kraus expression for that
second term.
\end{itemize}

Two quantum channels can lead to the same average-output purity. For instance,
consider the completely dephasing channel:
\begin{equation}
%\label{eq:comletely:dephasing:channel}
\mcE(\rho) = \frac{1}{2}\qty(\rho + \sigma_z \rho \sigma_z)
\end{equation}
and a depolarizing channel:
\begin{equation}
\mcE(\rho) = (1 - p) \rho + p \dsone,\quad p = \frac{1}{3}\qty(\sqrt{15} - 3).
\end{equation}
The average output-state purity for both channels is 
$\mathbb{E}_\psi \big[ \Tr(\mathcal{E}(\ket{\psi}\bra{\psi})^2) \big] = 2/3$.
Hence, evidencing that the average output-state purity may not be able to 
resolve differences.

Another example is the amplitude-damping channel, which maps the Bloch sphere 
onto a point on its surface. The average output-state purity for this process 
remains one, akin to any unitary transformation. In other words, unitary and 
non-unitary processes cannot be distinguished based solely on average 
output-state purity.

\jafix{Though, since the average output-state purity has a single global minimum, it 
is uniquely identified as a completely depolarizing channel.}{Si incluyo esto hay que argumentar más}

Other relevant example can be the amplitude-damping-like channel mapping 
to Bloch sphere onto $(0,0,0.01)$ is indistuinguishable from a dephasing-like
channel mapping the Bloch sphere onto a line $(0,0,0.017234z)$.

\clearpage
